\article{论自守表示的概念：前文的补充}{On the notion of an automorphic representation. A supplement to the preceding paper}{R. P. Langlands}{1}{203--207}

局部域（local field）上约化群（reductive group）的不可约表示（irreducible representation），可以由抛物子群（parabolic subgroup）的 Levi 因子的平方可积表示（square-integrable representation）出发，通过诱导（induction）并取子商（subquotient）得到 [2]、[4]。在整体域（global field）$F$ 上，类似地，从尖点表示（cuspidal representation）出发，可以得到所有自守表示（automorphic representation）；这里尖点表示对应于局部情形的平方可积表示。

设 $P$ 是 $G$ 的抛物子群，Levi 因子为 $M$，$\sigma=\bigotimes\sigma_v$ 是 $M(\A)$ 的一个尖点表示。那么 $\Ind\sigma=\bigotimes_v\Ind\sigma_v$ 是 $G(\A)$ 的表示，它未必不可约，甚至未必具有有限合成列（composition series）。按照惯例，将它的不可约子商称为成分（constituent）。

对几乎所有的 $v$，$\Ind\sigma_v$ 恰有一个包含 $G(O_v)$ 平凡表示（trivial representation）的成分 $\pi_v^\circ$。若 $\Ind\sigma_v$ 作用在 $X_v$ 上，则可按如下方式构造 $\pi_v^\circ$：取 $X_v$ 中包含所有非零 $G(O_v)$ 不变向量的最小 $G(F_v)$ 不变子空间 $V_v$，再取 $V_v$ 中不含这种向量的最大 $G(F_v)$ 不变子空间 $U_v$，最后让 $G(F_v)$ 作用在 $V_v/U_v$ 上。

\begin{originaltheorem}{引理 1（Lemma 1）}
$\Ind\sigma$ 的成分正是表示 $\pi=\bigotimes\pi_v$，其中 $\pi_v$ 是 $\Ind\sigma_v$ 的成分，并且对几乎所有的 $v$ 有 $\pi_v=\pi_v^\circ$。
\end{originaltheorem}
\begin{proof}
这样构造的表示显然都是成分。反过来，设成分 $\pi$ 作用在 $V/U$ 上，其中 $0\subseteq U\subseteq V\subseteq X=\bigotimes X_v$。回顾张量积（tensor product）的构造：选取有限的素位集 $S_0$，并对每个不在 $S_0$ 中的 $v$ 选取一个非零且被 $G(O_v)$ 固定的向量 $x_v^\circ$。可选取包含 $S_0$ 的有限素位集 $S$ 和 $X_S=\bigotimes_{v\in S}X_v$ 中的向量 $x_S$，使
\[
x=x_S\otimes\Bigl(\bigotimes_{v\notin S}x_v^\circ\Bigr)
\]
属于 $V$ 而不属于 $U$。

令 $V_S$ 为 $X_S$ 中包含 $x_S$ 且在 $G_S=\prod_{v\in S}G(F_v)$ 下不变的最小子空间。显然有满射
\[
V_S\otimes\Bigl(\bigotimes_{v\notin S}V_v\Bigr)\longrightarrow V/U;
\]
若 $v_0\notin S$，它的核包含
\[
V_S\otimes U_{v_0}\otimes\Bigl(\bigotimes_{v\notin S\cup\{v_0\}}V_v\Bigr).
\]
于是得到满射 $V_S\otimes(\bigotimes_{v\notin S}V_v/U_v)\longrightarrow V/U$，其核形如 $U_S\otimes(\bigotimes_{v\notin S}V_v/U_v)$，其中 $U_S\subseteq V_S$。$G_S$ 在 $V_S/U_S$ 上的表示不可约；又因为 $\Ind\sigma_v$ 有有限合成列，它形如 $\bigotimes_{v\in S}\pi_v$，其中 $\pi_v$ 是 $\Ind\sigma_v$ 的成分。因此 $\pi=\bigotimes\pi_v$，且当 $v\notin S$ 时 $\pi_v=\pi_v^\circ$。
\end{proof}

本文要证明以下命题。
\begin{originaltheorem}{命题 2（Proposition 2）}
$G(\A)$ 的表示 $\pi$ 是自守表示，当且仅当存在某个 $P$ 和某个 $\sigma$，使 $\pi$ 是 $\Ind\sigma$ 的一个成分。
\end{originaltheorem}

证明 $\Ind\sigma$ 的每个成分都是自守表示，需要运用 Eisenstein 级数（Eisenstein series）理论，而这一理论只在整体域 $F$ 的特征为零时得到了充分发展 [3]。不过可以期待，正特征整体域上的类似理论不久便会出现，因此作者在命题中没有特意限定特征为零。此外，证明每个自守表示都是某个 $\Ind\sigma$ 的成分，则可以不依赖 Eisenstein 级数理论。

先给出几个简单的引理。
\begin{originaltheorem}{引理 3（Lemma 3）}
设 $Z$ 是 $G$ 的中心。那么自守形式（automorphic form）是 $Z(\A)$ 有限的（$Z(\A)$-finite）。
\end{originaltheorem}
这一点已在 [1] 中证明。

\begin{originaltheorem}{引理 4（Lemma 4）}
设 $K$ 是 $G(\A)$ 的极大紧子群（maximal compact subgroup），$\varphi$ 是关于 $K$ 的自守形式。设 $P$ 是 $G$ 的抛物子群。选取 $g\in G(\A)$，并令 $K'$ 为 $M(\A)$ 的一个极大紧子群，包含 $gKg^{-1}\cap P(\A)$ 在 $M(\A)$ 上的投影。则
\[
\varphi_P(m;g)=\int_{N(F)\backslash N(\A)}\varphi(nmg)\,dn
\]
是 $M(\A)$ 上关于 $K'$ 的自守形式。
\end{originaltheorem}
增长条件（growth conditions）显然得以保留，且 $\varphi_P(\,\cdot\,;g)$ 光滑（smooth）并且是 $K'$ 有限的。它生成 $M$ 的局部 Hecke 代数（Hecke algebra）的容许表示（admissible representation），这一点可由 [2]、[4] 中的定理推出。

若对任意尖点形式 $\psi$ 和 $G(\A)$ 中的紧集 $\Omega$，都有
\[
\int_{\Omega G_{\mathrm{der}}(\A)}\varphi(g)\overline{\psi(g)}\,dg=0,
\]
就说 $\varphi$ 与尖点形式正交（orthogonal）。若 $P$ 是抛物子群，而对所有 $g$，$\varphi_P(\,\cdot\,;g)$ 都与 $M(\A)$ 上的尖点形式正交，则记为 $\varphi\perp P$。回顾如下简单引理 [3]。
\begin{originaltheorem}{引理 5（Lemma 5）}
若对所有 $P$ 都有 $\varphi\perp P$，则 $\varphi$ 为零。
\end{originaltheorem}

现在证明任意自守表示 $\pi$ 都是某个 $\Ind\sigma$ 的成分。可将 $\pi$ 实现在 $V/U$ 上，其中 $U,V$ 为自守形式空间 $\mathcal A$ 的子空间，且 $V$ 由单个自守形式 $\varphi$ 生成。在抛物子群的共轭类（conjugacy class）上规定全序，使 $\{P\}<\{P'\}$ 蕴含 $\operatorname{rank}P\leq\operatorname{rank}P'$，而 $\operatorname{rank}P<\operatorname{rank}P'$ 蕴含 $\{P\}<\{P'\}$。给定 $\varphi$，以 $\{P_\varphi\}$ 表示满足如下条件的最小 $\{P\}$：$\{P\}<\{P'\}$ 蕴含 $\varphi\perp P'$。在所有可用于生成 $\pi$ 的 $\varphi$ 中，选取使 $\{P\}=\{P_\varphi\}$ 最小的一个。对 $\psi\in\mathcal A$，令 $\psi_P(g)=\psi_P(1,g)$，并考虑 $V$ 上的映射 $\psi\mapsto\psi_P$。如果 $U$ 和 $V$ 的像相同，便可以把 $\pi$ 实现为此映射核的一个成分。这与 $\varphi$ 的选取方式矛盾。因此，若 $U_P,V_P$ 分别是 $U,V$ 的像，便可以在商空间 $V_P/U_P$ 中实现 $\pi$。

令 $\mathcal A_P^\circ$ 为 $N(\A)P(F)\backslash G(\A)$ 上满足下列条件的光滑函数 $\psi$ 的空间：
\begin{enumerate}
\item[(a)] $\psi$ 是 $K$ 有限的；
\item[(b)] 对每个 $g$，函数 $m\mapsto\psi(m,g)=\psi(mg)$ 是自守且尖点的。
\end{enumerate}
于是 $V_P\subseteq\mathcal A_P^\circ$。下文略去下标 $P$，仍将 $\pi$ 实现在 $V/U$ 上，其中 $U\subseteq V\subseteq\mathcal A_P^\circ$。假定 $V$ 由单个函数 $\varphi$ 生成。

\begin{originaltheorem}{引理 6（Lemma 6）}
设 $A$ 为 $M$ 的中心。可选取 $\varphi,V$，使得存在 $A(\A)$ 的特征（character）$\chi$，并且对所有 $g\in G(\A)$ 和 $a\in A(\A)$，有 $\varphi(ag)=\chi(a)\varphi(g)$。
\end{originaltheorem}
\begin{proof}
由于 $P(\A)\backslash G(\A)/K$ 是有限集，引理 3 蕴含 $\mathcal A_P^\circ$ 中任意 $\varphi$ 都是 $A(\A)$ 有限的。选取 $V$ 及其生成元 $\varphi$，使 $\{l(a)\varphi\mid a\in A(\A)\}$ 张成的空间 $Y$ 维数最小。这里 $l(a)$ 表示由 $a$ 给出的左平移。若维数为一，引理成立。否则存在 $a\in A(\A)$ 和 $\alpha\in\C$，使 $0<\dim(l(a)-\alpha)Y<\dim Y$。

分两种情形：$(l(a)-\alpha)U=(l(a)-\alpha)V$，或者两者不相等。后一情形下，可以用 $(l(a)-\alpha)\varphi$ 代替 $\varphi$，这与维数的最小性矛盾。前一情形下，可以将 $\pi$ 实现为 $V$ 中 $l(a)-\alpha$ 的核的一个子商。

接下来，在 $A(\A)$ 中选取格（lattice）$B$，使 $BA(F)$ 为闭集且 $BA(F)\backslash A(\A)$ 紧。在所有使 $Y$ 的维数尽可能小的 $\varphi,V$ 中，选取使 $B$ 的子群
\[
\{b\in B\mid l(b)\varphi=\beta\varphi,\ \beta\in\C\}
\]
秩最大的 $\varphi$。由上一段的论证可知，这个子群必须与 $B$ 同秩。由于 $\varphi$ 在 $A(F)$ 下不变，而 $BA(F)\backslash A(\A)$ 紧，可知此时 $Y$ 必须为一维，引理得证。
\end{proof}

选取这样的 $\varphi,V$，令 $\nu$ 为 $M(\A)$ 的正特征，满足
\[
\nu(a)=|\chi(a)|,\qquad a\in A(\A),
\]
考虑 Hilbert 空间 $L_2^\circ=L_2^\circ(M(F)\backslash M(\A),\chi)$，它由 $M(\Q)\backslash M(\A)$ 上满足下列条件的可测函数（measurable function）$\psi$ 组成：\footnote{\textbf{译者注（本次校订）：}本段及下方积分中的 $M(\Q)$ 应读作 $M(F)$，与全文的整体域记号一致。此外，$L_2^\circ$ 表示尖点部分，故还须要求 $\psi$ 为尖点函数；仅有下面两条条件，所得整个 $L^2$ 空间一般含连续谱，不能直接作下文的离散直和分解。}
\begin{enumerate}
\item[(i)] 对所有 $m\in M(\A)$ 和 $a\in A(\A)$，有 $\psi(ag)=\chi(a)\psi(g)$；
\item[(ii)] $\displaystyle\int_{A(\A)M(\Q)\backslash M(\A)}\nu^{-2}(m)|\psi(m)|^2\,dm<\infty$。
\end{enumerate}
$L_2^\circ$ 是不可约不变子空间的直和（direct sum）；若 $\psi\in V$，则对所有 $g\in G(\A)$，$m\mapsto\psi(m,g)$ 属于 $L_2^\circ$。选取 $L_2^\circ$ 的一个不可约分量 $H$，使某个 $g\in G(\A)$ 对应的 $\varphi(\,\cdot\,,g)$ 在 $H$ 上的投影非零。

对每个 $\psi\in V$，定义 $\psi'(\,\cdot\,,g)$ 为 $\psi(\,\cdot\,,g)$ 在 $H$ 上的投影。容易看出，对所有 $m_1\in M(\A)$，$\psi'(mm_1,g)=\psi'(m,m_1g)$。因此可令 $\psi'(g)=\psi'(1,g)$，定义函数 $\psi'(g)$。若 $V'=\{\psi'\mid\psi\in V\}$，则可将 $\pi$ 实现为 $V'$ 的一个商。若 $\delta^2$ 是 $M(\A)$ 在 $N(\A)$ 上的模函数（modular function），$\sigma$ 是 $M(\A)$ 在 $H$ 上的表示，则 $V'$ 包含在 $\Ind\delta^{-1}\sigma$ 的空间中。

逆命题的证明要用到由尖点形式构造的 Eisenstein 级数及其解析延拓（analytic continuation）。设 $\pi$ 是相对于 $G(\A)$ 的某个极大紧子群 $K$ 定义的整体 Hecke 代数 $\mathcal H$ 的一个表示。选取 $\pi$ 中包含的一个 $K$ 的不可约表示 $\kappa$。若 $E_\kappa$ 是由 $\kappa$ 定义的幂等元（idempotent），令 $\mathcal H_\kappa=E_\kappa\mathcal HE_\kappa$，并令 $\pi_\kappa$ 为 $\mathcal H_\kappa$ 在 $\pi$ 的 $\kappa$ 等型子空间（isotypical subspace）上的不可约表示。要证明 $\pi$ 是自守表示，只需证明 $\pi_\kappa$ 是 $\mathcal H_\kappa$ 在 $\kappa$ 型自守形式空间上的表示的一个成分。为简化记号，以下以 $\pi$ 表示 $\pi_\kappa$，以 $\mathcal H$ 表示 $\mathcal H_\kappa$。

固定 $P$ 及 $M(\A)$ 的尖点表示 $\sigma$。令 $L$ 为 $M$ 的 $F$ 有理特征（rational character）所成的格，$L_\C=L\otimes\C$。$L_\C$ 的每个元素 $\mu$ 定义 $M(\A)$ 的一个特征 $\xi_\mu$。令 $I_\mu$ 为 $\Ind\xi_\mu\sigma$ 的 $\kappa$ 等型子空间，$I=I_0$。要证明：若 $\pi$ 是 $I$ 上表示的一个成分，那么 $\pi$ 是 $\mathcal H$ 在 $\kappa$ 型自守形式空间上的表示的一个成分。

若 $\{g_i\}$ 为 $P(\A)\backslash G(\A)/K$ 的一组陪集代表，可由映射 $\varphi\mapsto\varphi_\mu$ 将 $I_\mu$ 与 $I$ 认同，其中
\[
\varphi_\mu(nmg_i k)=\xi_\mu(m)\varphi(nmg_i k).
\]
换言之，这给出了 $L_\C$ 上丛（bundle）$\{I_\mu\}$ 的一个平凡化（trivialization），于是可以谈论全纯截面（holomorphic section），或在 $\mu=0$ 处消失到某个阶数的截面。这些概念不依赖于 $g_i$ 的选择，尽管平凡化本身依赖于该选择。

存在 $\mu=0$ 的邻域 $U$ 以及与 $U$ 相交的有限个超平面，使得对 $U$ 中这些超平面之外的 $\mu$，$\varphi\in I_\mu$ 的 Eisenstein 级数 $E(\varphi)$ 有定义。将 $E$ 乘以适当的线性函数的乘积，可以消去这些超平面上的极点，使它在整个 $U$ 上都有定义。以下只使用这个修正后的函数，仍记作 $E$。它取值于自守形式空间，因而 $E(\varphi)$ 是 $G(\A)$ 上的函数 $g\mapsto E(g,\varphi)$。若 $h\in\mathcal H$，且 $\rho_\mu$ 为 $\Ind\xi_\mu\sigma$，则满足
\[
E(\rho_\mu(h)\varphi)=r(h)E(\varphi).
\]
此外，若 $\varphi_\mu$ 是 $\{I_\mu\}$ 在 $0$ 的邻域内的全纯截面，则对每个 $g$，$E(g,\varphi_\mu)$ 关于 $\mu$ 全纯，逐点对 $E(\varphi_\mu)$ 求导，所得函数仍属于 $\mathcal A$。

令 $I_r$ 为 $I$ 的全纯截面在 $\mu=0$ 处的 $r$ 阶芽（germ）的空间。那么 $\varphi_\mu\mapsto\rho_\mu(h)\varphi_\mu$ 定义了 $\mathcal H$ 在 $I_r$ 上的作用。若 $s\leq r$，自然映射 $I_r\to I_s$ 是 $\mathcal H$ 同态（homomorphism）；记其核为 $I_r^s$。显然 $I_0=I$。选取 $L_\C$ 上线性形式的一组基，可以考虑取值于 $\mathcal A$ 的 $\kappa$ 等型子空间的幂级数 $\sum_{|\alpha|\leq r}\mu^\alpha\psi_\alpha$。$\mathcal H$ 通过右平移作用在这一空间上，所得表示是 $\kappa$ 等型自守形式上的表示的直和。此外，$\varphi_\mu\mapsto E(\varphi_\mu)$ 定义了从 $I_r$ 到这一空间的 $\mathcal H$ 同态 $\lambda$。结合 Jordan--Hölder 定理与下面的引理，就能完成命题的证明。

\begin{originaltheorem}{引理 7（Lemma 7）}
当 $r$ 充分大时，$\lambda$ 的核包含于 $I_r^0$。
\end{originaltheorem}
由于 $P,\sigma$ 已经固定，可以用沿与 $P$ 相伴的抛物子群的常数项（constant term）之和代替 $E$，并相应修正 $\lambda$。所有这些常数项恒等于零，当且仅当 $E$ 本身恒等于零。若 $Q_1,\ldots,Q_m$ 是与 $P$ 相伴的抛物子群的各个类的一组代表，令 $E_i(\varphi)$ 为沿 $Q_i$ 的常数项。可假定 $M$ 是每个 $Q_i$ 的一个 Levi 因子。对 $m\in M(\A)$，由 $\xi_\mu(m)=e^{\langle\mu,\nu(m)\rangle}$ 定义 $\nu(m)$，于是 $\nu(m)$ 属于 $L_\R$ 的对偶。若 $\varphi\in I_\mu$，函数 $E_i(\varphi)$ 形如
\[
E_i(nmg_jk,\varphi)=\sum_{\alpha=1}^{a}\sum_{\beta=1}^{b}
 p_\alpha(\nu(m))\xi_{\nu_\beta(\mu)}(m)\psi_{\alpha\beta}(m,k).
\]
这里 $\psi_{\alpha\beta}$ 属于一个不依赖于 $\mu,g_j$ 的有限维空间；$\nu_\beta$（$1\leq\beta\leq b$）是 $\mu$ 的全纯函数；$\{p_\alpha\}$ 是某个次数的多项式的一组基。这种表达式未必唯一。下一个引理表明，可以缩小开集 $U$，再找到 $G(\A)$ 中的有限集 $h_1,\ldots,h_n$，使得对 $\mu\in U$、$\varphi_\mu\in I_\mu$，$E(\varphi_\mu)=0$ 当且仅当所有数 $E_i(h_j,\varphi_\mu)$（$1\leq i\leq m$，$1\leq j\leq n$）均为零。

\begin{originaltheorem}{引理 8（Lemma 8）}
设 $U$ 是 $\C^l$ 中 $0$ 的一个邻域，$\nu_1,\ldots,\nu_k$ 是 $U$ 上的全纯函数，$p_1,\ldots,p_a$ 是某个给定次数的多项式的一组基。那么存在包含于 $U$ 的 $0$ 的邻域 $V$ 和 $\C^l$ 中有限集 $y_1,\ldots,y_b$，使得当 $\mu\in V$ 时，等式
\[
\sum p_i(y)e^{\nu_j(\mu)\cdot y}=0\tag{*}
\]
对所有 $y$ 成立，当且仅当它对 $y=y_1,\ldots,y_b$ 成立。
\end{originaltheorem}
为证明此引理，只需注意，由方程 (*)（$y\in\C^l$）定义的 $U$ 的解析子集（analytic subset），在 $0$ 的一个邻域内由其中有限个方程定义。

因此可以把 $E$ 看作 $U$ 上的一个函数，其值为从有限维空间 $I$ 到 $\C^{mn}$ 的线性变换。由 Eisenstein 级数理论可知，当 $\mu$ 属于 $U$ 的一个开子集时，$E_\mu$ 是单射（injective）。因此，只需再证明下面的引理。

\begin{originaltheorem}{引理 9（Lemma 9）}
设 $E$ 是 $\C^l$ 中 $0$ 的邻域 $U$ 上取值于 $\Hom(I,J)$ 的全纯函数，其中 $I,J$ 为有限维空间，并设 $E_\mu$ 在 $U$ 的一个开子集上为单射。那么存在整数 $r$，使得若 $\varphi_\mu$ 在 $\mu=0$ 附近解析，而 $E_\mu\varphi_\mu$ 的 Taylor 级数消失到 $r$ 阶，则 $\varphi_0=0$。
\end{originaltheorem}
\begin{proof}
投影到 $J$ 的一个商空间后，可以假定 $\dim I=\dim J$，甚至 $I=J$。设 $\det E_\mu$ 的幂级数展开中第一个非零项的次数为 $s$。下面证明 $r=s+1$ 即可。只需对 $l=1$ 证明，因为可以限制到一条直线上，使 $E_\mu$ 的首项仍为 $s$ 次。此时在 $E$ 左右分别乘以适当的可逆全纯矩阵，可将其化为对角阵，对角元为 $z^\alpha$，$0\leq\alpha\leq s$。于是结论显然。
\end{proof}
最后，感谢 P. Deligne，他指出了本文初稿中的一个错误。

\begin{thebibliography}{4}
\bibitem[1]{c13-1} A. Borel and H. Jacquet, \emph{Automorphic forms and automorphic representations}, 本论文集，第1册，189--202页。
\bibitem[2]{c13-2} P. Cartier, \emph{Representations of $p$-adic groups: A survey}, 本论文集，第1册，111--155页。
\bibitem[3]{c13-3} R. P. Langlands, \emph{On the functional equations satisfied by Eisenstein series}, Lecture Notes in Math., vol. 544, Springer, New York, 1976.
\bibitem[4]{c13-4} N. Wallach, \emph{Representations of reductive Lie groups}, 本论文集，第1册，71--86页。
\end{thebibliography}
\noindent 高等研究院（The Institute for Advanced Study）。

\noindent\footnotesize 原文 AMS（MOS）学科分类（1970）：主分类 30A58、20E35、10D20。版权 © 1979 American Mathematical Society。\normalsize
